\chapter{结论与展望}

本文围绕高速公路团雾监测这一典型低能见度智能交通问题，提出并实现了一种基于深度估计与多任务学习的统一监测方法。针对真实团雾样本稀缺、浓度标签难以获取、固定监控视频场景与自动驾驶数据差异明显等问题，本文选择以 UA-DETRAC 清晰交通数据为基础，通过 MiDaS\_small 预计算深度缓存，构建基于大气散射模型的在线造雾机制，在训练阶段动态生成 clear、uniform、patchy 三类天气样本，并进一步设计统一多任务模型同时完成车辆检测、雾类型分类和 beta 浓度回归。与单一任务研究相比，本文方法更贴近高速公路监控场景下的真实业务需求。

在工程实现方面，本文完成了从数据准备、训练、验证、视频推理到导出与量化路径的完整闭环。当前工作区已经能够稳定读取原始数据、补齐深度缓存、执行联合训练并输出模型权重；推理脚本能够在视频中显示雾类别、beta 数值、车辆检测结果和动态阈值控制效果；混合推理脚本则进一步解决了早期统一模型在真实雾视频上车辆检测能力不足的问题，并构建了带中文标题、刻度条和分级标签的雾浓度分布图小地图，从而显著提升了系统演示的完整性与可解释性。

从研究结论看，本文工作至少说明了以下几点。第一，在高速公路固定监控视频场景中，利用单目深度估计和在线物理造雾构造多任务训练监督是可行的。第二，将车辆检测、雾分类和雾浓度回归整合为统一多任务模型具有合理性，并已在工程上打通训练与推理主链。第三，基于 beta 的时间平滑和检测阈值联动，为天气感知参与业务决策提供了初步可执行路径。第四，即使统一模型尚在持续优化，通过混合推理也可以较快形成可视化系统原型，为论文验证和后续工程迭代争取时间。

当然，本文仍处于毕业设计原型系统阶段，尚有若干值得深入推进的方向。未来工作可从以下几个方面展开。

首先，应继续强化真实雾域适应能力。一种可行做法是在当前在线造雾训练框架基础上，引入少量真实雾视频做有监督或半监督微调，以缩小合成雾与真实雾之间的域差距。

其次，可将当前的全局 beta 回归扩展为空间雾浓度预测。若能设计一个密集回归分支输出逐像素或分块级雾浓度图，则右下角的小地图就能够从“解释增强近似图”升级为真正的模型预测结果。

再次，可引入更直接的能见度标定数据。如果后续能够获得路侧气象站、能见度仪、激光雷达或人工等级标注数据，就可以将 beta 与真实能见度建立更明确对应关系，增强结果的物理意义。

最后，在部署层面，本文已预留 QAT、INT8 和 ONNX 路径，未来可以进一步围绕 Jetson 或路侧边缘设备开展推理速度、功耗、时延和稳定性测试，使该系统真正走向“可部署、可维护、可联动”的交通气象智能监测原型。

若进一步结合陕西省交通运输行业现状，未来工作还可以围绕三个更具体的方向推进。第一，面向陕西交控集团和省级路网运行部门的既有监测预警项目，探索将本文的视频雾风险识别能力以插件方式接入现有风险点工程与值班调度流程；第二，围绕西安绕城、新机场线、连霍高速关中段和陕南山区桥隧通道，分别构建“高流量平原场景”和“高风险山地场景”两类验证样板，比较不同区域的雾特征差异与部署重点；第三，结合陕西交通部门已经形成的气象预警转发、路网周报和省际应急演练机制，进一步把模型输出从单纯的视频可视化结果提升为可直接进入限速诱导、情报板发布和会商研判的结构化信息。若这些工作能够继续推进，本文项目在陕西场景中的应用价值将不仅体现为一套毕业设计原型，而有望成为面向实际行业系统的小型智能模块。

总的来说，本文在高速公路团雾监测这一具有明显现实意义的问题上，完成了一套从理论方法到工程系统的完整探索。虽然仍有大量细节有待完善，但其方法框架、代码实现和可视化闭环已经具备较强的继续扩展价值，可为毕业设计终稿、后续论文发表以及实际道路监测应用提供坚实基础。
